\section{Data and Protocol}
\label{sec:data-protocol}

\subsection{Data}
\label{sec:data}

Experiments use the Desert Knowledge Australia Solar Centre (DKASC)
dataset at Alice Springs (23.767$^\circ$S, 133.867$^\circ$E, 558 m), an
open dataset of ground-mounted PV installations with per-array power
measurement and a co-located weather station \cite{dkasolarcentre},
also used by six of the papers in our survey
\cite{zhou2024cnnlstmattnbayes, hou2024vmdwoalstm,
ye2026distributedcnnlstm, hussain2022hybridgrucnn,
alharkan2023dsclanet, guo2024qrkddn}: every number in this paper is
independently checkable against it.
Three fixed-mount arrays are used, chosen for differing module
technology at comparable scale: array 11 (BP Solar, polycrystalline
silicon, 5.0 kW), array 12 (BP Solar, monocrystalline silicon, 5.1
kW), and array 17 (Sanyo, heterojunction, 6.3 kW), all mounted at
20$^\circ$ tilt. The three arrays share one weather station and are
not independent replications of one another.

Splits are chronological whole calendar years: training 2011-2013,
validation 2014, test 2015, the last touched once, at the end.
Training begins in 2011 because array 17 was commissioned in March
2010, and all three arrays share this training window so that
cross-array comparisons are not confounded by differing training
length.

Evaluation is restricted to hours with solar elevation above
10$^\circ$, stricter than the 85$^\circ$ zenith filter Yang et
al.~\cite{yang2020verification} describe as typical practice, applied
for the reason they give: at low solar elevation the clear-sky index
becomes unstable relative to a small denominator. One further
exclusion is applied: array 17 was switched off for a documented
maintenance outage from 5 to 9 June 2015, excluded from evaluation for
that array only.

\subsection{The Excluded Array}
\label{sec:data-excluded}

A fourth array, array 07 (First Solar, cadmium telluride, 7.0 kW), was
initially included and subsequently excluded, and is reported rather
than omitted, because the mechanism by which it was nearly retained is
itself evidence for this harness. A conventional completeness audit
passed array 07 for 2014 with 99.99 percent coverage and 0.00 percent
missing values in the power channel; the array had in fact produced no
power from March to September that year, with 48.41 percent of its
daylight hours recording exactly zero output. The audit did not detect
this because the logger continued to record faithfully, and what it
recorded was zero: zero is not missing, and a completeness-only check
is structurally blind to a healthy sensor reporting a dead array. A
second audit, testing for near-zero rather than exactly-missing output,
is what caught it, and both audits are included in the anonymised
repository accompanying this paper.

\subsection{Protocol}
\label{sec:protocol}

We state here only the rules bearing on this paper's claim; the full
evaluation protocol, including feature construction and significance
testing, is available in the anonymised repository accompanying this
paper.

Splits are never shuffled, and all scalers and feature statistics used
by the models are fitted on training data only.

Three reference forecasts are used. Smart persistence forecasts
$\hat{P}(t) = k_p(t-h) \times P_{cs}(t)$, persisting the clear-sky
index from the issue time. Climatology forecasts $\hat{P}(t) =
\bar{k}_p(\text{month}, \text{hour}) \times P_{cs}(t)$, with the mean
clear-sky index per calendar month and hour of day estimated on
training data only. The convex reference is $w \times$ persistence $+
(1-w) \times$ climatology, with $w$ selected by grid search to
minimise root-mean-square error on the validation split, fitted
independently for each array, horizon, seed and regime. Because $w$ is
chosen to minimise error on that split, validation skill against the
convex reference is in-sample for the reference; $w$ is never refit on
test, so the held-out test year is the clean, out-of-sample
measurement, and Section~\ref{sec:rq1-reference} reports it first.
Following Yang et al.~\cite{yang2020verification}, the convex
reference is the headline metric; skill against plain persistence is
reported alongside because the gap between the two is a result, not a
robustness check.

The skill score is $s = 1 - \text{RMSE}_{\text{forecast}} /
\text{RMSE}_{\text{reference}}$, computed against both references over
the same forecasts and the same evaluation sample.
